\documentclass[../main.tex]{subfiles}

\begin{document}

\chapter{Lebesgue 积分}\label{chapter:ch04-Lebesgue积分}

我们现在转向第一部分中感兴趣的主要目标——Lebesgue 积分. 我们用四个步骤定义该积分. 首先对在有限测度集上的简单函数定义该积分. 接着对有限测度集上的有界可测函数定义, 这通过用简单函数根据 $f$ 的上逼近与下逼近的积分来定义. 我们定义 $E$ 上的一般的非负可测函数的积分为在有限测度集外为零的有界可测函数的下逼近的积分的上确界, 这样的函数的积分是非负的, 但可能是无限的. 最后, 一般的可测函数称为在 $E$ 上可积, 若 $\int_{E}|f|<\infty$. 我们证明可积函数的线性组合是可积的. 在可积函数类上, Lebesgue 积分是一个单调的、线性的泛函. Lebesgue 积分的一个主要优点, 除了可积函数类的拓展外, 是以下十分普通的准则的应用, 即确保若存在可积函数序列 $\{f_{n}\}$ 在 $E$ 上几乎处处逐点收敛于 $f$, 则
\begin{align*}
\lim_{n\to\infty}\int_{E}f_{n}=\int_{E}\left[\lim_{n\to\infty}f_{n}\right]\equiv\int_{E}f.
\end{align*}
我们称之为积分号下取极限. 基于 Egoroff 定理, 它是 Lebesgue 测度的可数可加性的推论, 我们证明了四个提供这种取极限的合理性准则的定理: 有界收敛定理, 单调收敛定理, Lebesgue 控制收敛定理以及 Vitali 收敛定理.

\subfile{Chapter-04/section-01.tex}

\subfile{Chapter-04/section-02.tex}

\subfile{Chapter-04/section-03.tex}

\subfile{Chapter-04/section-04.tex}

\subfile{Chapter-04/section-05.tex}

\subfile{Chapter-04/section-06.tex}

\end{document}